\section{Discussion}

\subsection{Imbalance correction is architecture-dependent}

Focal loss with class-balanced sampling is not a neutral intervention. It helps the transformer and hurts both CNNs. A benchmark that applied the correction to one architecture and generalized the result would have reached the wrong conclusion about the others.

The effect has a plausible mechanism. Focal loss changes the gradient distribution, and class-balanced sampling changes the effective class prior. A transformer with weaker inductive bias may benefit from the harder examples and the reweighted prior, while a CNN with strong local features may overfit the minority samples or lose calibration. The experiment cannot separate those mechanisms. We report the effect as an empirical interaction, not a causal explanation.

\subsection{The ranking flip and the stable AUC}

The sensitivity ranking reverses, but the AUC ranking does not. EfficientNet-B0 remains first and ViT-S/16 remains last in both the baseline and corrected conditions. The reversal is therefore specific to the thresholded operating point.

This matters for how model comparisons are reported. A single threshold at 0.5, or a single sensitivity number, can make an architecture look better or worse depending on the correction. Reporting sensitivity at a fixed specificity, AUC, and calibration together gives a more complete picture. We leave a full calibration analysis to future work.

\subsection{The fairness signal}

The dark stratum has the lowest AUC and the highest false positive rate in all nine configurations. The direction is consistent even though the sample is small. The most likely interpretation is that the model's ranking quality degrades on darker images, which produces more false positives at the fixed specificity threshold.

The dark-stratum sensitivity is underpowered. Six melanoma cases cannot support a precise estimate, and the bootstrap intervals span the full range. We therefore treat the AUC and false positive rate as the more stable signals and the sensitivity estimate as inconclusive.

The light stratum is often harder than the medium stratum. This is consistent with a contrast effect: a pale lesion on pale skin has a lower boundary contrast, and the model may attend to the surrounding skin. Skin tone is not the only source of difficulty. Lesion contrast, hair, rulers, and illumination can all change the model's attention.

\subsection{The data bottleneck}

The external dermoscopic datasets with Fitzpatrick labels, HIBA and PROVe-AI, contain Fitzpatrick I-IV only. Pipsqueak, the only dermoscopic dark-skin set, has 16 core images. The clinical datasets with Fitzpatrick V-VI, such as DDI, SCIN, and PASSION, differ from HAM10000 in imaging modality.

This is the main obstacle to a strong dark-skin claim. This study can quantify the representation gap, report the consistent dark-stratum AUC and false-positive signal, and show the limitation of the available dermoscopic data. A precise dark-skin sensitivity estimate requires new dermoscopic data collection or a carefully controlled cross-modality study.

\subsection{What XAI adds}

The classification and fairness results describe what the model does. The XAI analysis tests whether it attends to the lesion or to the surrounding skin. If attribution mass outside the lesion is higher in the dark stratum and correlates with errors, attribution leakage becomes an early warning signal that can be measured before group-level sensitivity collapses. The protocol is defined in Section~3.9.

\subsection{Implications for model selection}

The choice of architecture depends on the correction. At the CE baseline, ResNet-50 has the highest sensitivity at 95\% specificity. Under focal loss with class-balanced sampling, ViT-S/16 is the best. A deployment decision that ignores the training condition can pick the wrong model.

For the fairness analysis, this means the evaluation should include the architecture-condition combinations, not a single global best model. The best model on aggregate sensitivity may not be the best model in the dark stratum, and the ranking may change again after the external validation.
